\documentclass[11pt,letterpaper]{article}
\usepackage[T1]{fontenc}
\usepackage[utf8]{inputenc}
\usepackage{lmodern,microtype}
% Explicitly load the standard Euler font map on minimal pdfTeX installations.
\ifdefined\pdfmapfile\pdfmapfile{+euler.map}\fi
\usepackage[margin=0.9in,headheight=15pt]{geometry}
\usepackage{amsmath,amssymb,amsthm,mathtools}
\usepackage{booktabs,longtable,array,tabularx}
\usepackage[dvipsnames]{xcolor}
\usepackage{enumitem,fancyhdr,needspace}
\usepackage{xurl}
\usepackage[colorlinks=true,linkcolor=MidnightBlue,urlcolor=MidnightBlue,citecolor=MidnightBlue]{hyperref}
\hypersetup{pdftitle={Finite Frobenius certificates for the theta exponent},pdfauthor={ChatGPT},pdfsubject={A source-based simplification proposal for the submitted RH and Dirichlet GRH development}}
\definecolor{ink}{HTML}{193B50}
\definecolor{pale}{HTML}{F0F5F7}
\pagestyle{fancy}
\fancyhf{}
\fancyhead[L]{\small\sffamily Finite Frobenius certificates}
\fancyhead[R]{\small\sffamily GRH source analysis}
\fancyfoot[C]{\thepage}
\renewcommand{\headrulewidth}{0.3pt}
\setlength{\parindent}{0pt}
\setlength{\parskip}{5.5pt plus 1pt}
\setlist{nosep,leftmargin=1.6em,topsep=5pt}
\setcounter{tocdepth}{1}
\newtheorem{theorem}{Theorem}[section]
\newtheorem{lemma}[theorem]{Lemma}
\newtheorem{corollary}[theorem]{Corollary}
\theoremstyle{remark}
\newtheorem{remark}[theorem]{Remark}
\newcommand{\OO}{\mathcal O}
\newcommand{\JJ}{\mathcal J}
\newcommand{\FF}{\mathsf F}
\newcommand{\dd}{\mathrm d}
\newcommand{\ord}{\operatorname{ord}}
\newcommand{\coef}[2]{[\, #1\,]#2}
\newcommand{\code}[1]{{\small\nolinkurl{#1}}}
\newcommand{\src}[1]{\hyperref[src:#1]{[S#1]}}
\newcommand{\repo}{https://github.com/VladimirReshetnikov/ProveIt}
\newcommand{\commit}{9050b04000647d588d0c78d945180abcd3b0bcf7}
\newcommand{\sourceurl}[1]{\repo/blob/\commit/Analysis/GeneralizedRiemannHypothesis/#1}
\newcommand{\sourcelink}[2]{\href{\sourceurl{#1}}{#2}}
\newcommand{\sectionlead}[1]{\needspace{6\baselineskip}\textbf{#1}}

\begin{document}
\thispagestyle{empty}
{\color{ink}\sffamily\bfseries\LARGE Finite Frobenius certificates\par}
\vspace{3pt}
{\color{ink}\sffamily\Large Replacing the theta exponential normal form\par}
\vspace{8pt}
{\large A simplification proposal for the submitted RH / Dirichlet-GRH development}

\vspace{7pt}
\textbf{Source snapshot:} \code{9050b04000647d588d0c78d945180abcd3b0bcf7}\par
\textbf{Report date:} 27 November 2026 (Pacific time)\par
\textbf{Deliverables:} mathematical replacement proof, checked standalone Lean core,
source audit, exact finite-ring regressions, and a staged integration plan.

\vspace{8pt}
\colorbox{pale}{\parbox{\dimexpr\linewidth-2\fboxsep\relax}{
\textbf{Recommendation.} Keep the raw integral exponent of the theta product.
Replace its complete expansion as
\(\sum_{e,j}\lambda_{e,j}\Gamma_e^{p^j}\)
by two smaller contracts: finite null-moment vanishing for divisibility, and a
first-derivative signature for inert-prime sharpness. Finite Frobenius orthogonality
proves the first contract; differentiating the raw parameter equation proves the
second. Both use inputs that the current source already establishes.
}}

\begin{abstract}
The current local argument determines an infinite, two-index coupling expansion,
turns it into an exact exponential normal form, and then uses only selected
consequences of that normal form. The replacement developed here avoids this
intermediate classification. In a finite jet algebra, Frobenius is nilpotent on
the augmentation ideal. The scaled formal-parameter equation can therefore be
inverted by a finite geometric sum. A weighted pairing of two Frobenius iterates
is a Frobenius iterate of a single moment. Even and odd iterates are consequently
orthogonal when the inert coupling moments vanish; the ramified case uses all
moments, including degree two. Differentiating the raw exponent separately recovers
the same pure-power gradient and unit leading coefficients used for sharpness.

The replacement preserves the existing local valuation bounds and their global
error budgets. It targets 752 physical lines in four source files, with another
84 helper lines subject to a dependency audit. These are affected regions, not
measured net deletions. A 314-line independent Lean file checks the finite algebraic
core against the repository's exact Lean and mathlib versions. The formal-group
adapters and full-project refactor remain to be implemented. The proposal is
distinct from the mechanisms and stated future directions in the 18 existing
proposal directories at the pinned commit.
\end{abstract}

\clearpage
\thispagestyle{plain}
\tableofcontents
\newpage

\section{What the proof is doing, and where to intervene}
\label{sec:architecture}

\subsection{The endpoint and the quantitative spine}

The exported theorem concerns all Dirichlet characters, with the usual exception
at the pole of the principal character. In the source, \code{New.grh} proves
nonvanishing for \(\Re s>1/2\); the functional equation then gives the line statement,
and the modulus-one case gives the Mathlib statement of RH. Here ``GRH'' is the
Dirichlet family specified by those exported declarations. \src{1}

The arithmetic object is the smoothed prime remainder
\begin{align}
 \lambda_u(n)&=\mathbf 1_{n\equiv1\pmod u}\Lambda(n)(1-\chi_d(n)),\notag\\
 R_u(x)&=\sum_{1\le n\le x}\lambda_u(n)\left(\frac1n-\frac1x\right)
       -\frac{\log x}{\varphi(u)}-c_u. \label{eq:remainder}
\end{align}
For admissible fixed \(Q,d\), the target is the one-sided estimate
\(R_Q(x)\le x^{-1/2+\varepsilon}\) eventually, for every \(\varepsilon>0\).
The guide and the final reduction explain how suitable quadratic twists and the
one-sided zero-detection criterion turn this estimate into nonvanishing. \src{2}\src{3}

The algebraic comparison has a particularly clear structure. Let \(d_T\) be the
target dimension, \(H\) its main height term, \(\mathcal R\) the total prime
remainder, and \(m\) the chain length. Write
\[
 L=fa(\log b)^4,\qquad I=10\log(1+2b^2)\,fa(\log b)^3.
\]
The three inputs are:
\begin{center}
\begin{tabularx}{\linewidth}{@{}l X@{}}
\toprule
Input & Exact role in the comparison\\
\midrule
A: short forms & Independent candidate images have squared lengths at most
\(\exp(H-m\mathcal R+L)\); their denominators have a controlled index.\\
B: enough forms & Choose \(n\) independent such images with
\((1-4/b^2)d_T\le n\le d_T\), using few enough translates that the index is at most \(d_TI\).\\
C: long forms & For that same family, squared covolume plus twice the index is at least \(d_T(H-3L)\).\\
\bottomrule
\end{tabularx}
\end{center}
Hadamard's inequality therefore gives
\[
 d_T(H-3L)\le n(H-m\mathcal R+L)+2d_TI.
\]
Since \(H+L\ge0\) and \(n\le d_T\), the main term cancels to give
\[
 n\,m\mathcal R\le d_T(4L+2I).
\]
Using \(n\ge d_T/2\) and the parameter bounds yields the source theorem
\begin{equation}
 m\mathcal R\le fa(\log b)^5. \label{eq:star}
\end{equation}
The precise normalization and the matching main term matter much more than a
small improvement in a numerical constant. The final scales have chain length
on the square-root scale and a very distant second cutoff; they convert
\eqref{eq:star} to the one-sided prime estimate. \src{3}\src{4}

\subsection{Why the proposed change is local}

The existing proposals already attack the main global opportunities: spanning,
rank/completion budgets, twist selection, signed normalization, torsion
interpolation, local labels, invariant averages, class counting, and orbital
exchange. Appendix~\ref{app:inventory} records the full exclusion inventory.
Repeating one of those reductions, or merely asking for a stronger determinant
estimate mentioned there, would not answer the request.

The remaining concrete opportunity identified here is an interface inside the
local derivative calculation. Finite-place divisibility makes the upper forms
short; matching sharpness also enters the lower determinant calculation. Thus a
replacement has to preserve \emph{both} consequences. An argument that keeps only
integrality would not justify retaining the global comparison.

This report analyzes the proof architecture and supplies a replacement for this
specific interface. It does not assert a fresh independent replay of the entire
GRH development. The source revision is fixed, the exact algebra checked here is
identified, and the remaining integration work is explicit.

\section{The stronger object currently constructed}
\label{sec:old}

At a nonsplit prime \(p\), let \(\OO\) be the valuation ring of
\(\mathbb C_p\), normalize \(v(p)=1\), and let \(\pi=\pi_p\) be the element
called \code{piPrimeInt d}. Let \(\varpi\) be the chosen scale and set
\[
 c=\frac{\varpi^{p+1}}p.
\]
The relevant couplings have the form
\begin{equation}
 \Gamma_e(z)=\sum_{x}z_x^{p^e}z_{\operatorname{pr}_e(x)}.
 \label{eq:couplings}
\end{equation}
The variables can be indexed by a place and a coordinate; the partner map changes
the place and keeps the coordinate. The allowed exponents are odd positive
integers at inert primes and all nonnegative integers at the ramified prime.

The present implementation introduces
\begin{equation}
 \operatorname{couplingSum}
   =\sum_{e\in S}\sum_{j\ge0}\lambda_{e,j}\Gamma_e^{p^j}
 \label{eq:normal}
\end{equation}
and proves an exact factorization
\begin{equation}
 F(\varpi z)=\widetilde W(z)
       \exp\!\left(c\operatorname{couplingSum}(z)\right),
 \qquad \widetilde W\in\OO[[z]]. \label{eq:exactnormal}
\end{equation}
It obtains this by first identifying the residue of the one-curve exponent,
expanding products of prime-power series, classifying representations of the
form \(p^a+p^b\), reindexing an infinite sum, summing the resolution factors,
and absorbing a divisible error into an integral exponential. Substitution,
scaling and differentiation are then proved for the resulting infinite normal
form. \src{5}\src{6}\src{7}\src{8}

The two downstream needs are smaller:
\begin{enumerate}
\item After a zero-constant integral substitution that kills relevant couplings
modulo \(\pi\), the raw exponent must vanish modulo \(\pi\) to the finite
degree being tested.
\item At inert primes, the first partial derivatives of the raw exponent must
have the prescribed pure prime-power monomials modulo the maximal ideal, with
unit coefficients.
\end{enumerate}
The proposal is to retain the exact factorization with the \emph{raw} exponent
and prove these two contracts directly. It does not posit an equality between
the raw exponent and a short sum of couplings. Such an equality is unnecessary.

\section{Finite Frobenius algebra}
\label{sec:algebra}

Let \(p\) be an odd prime and \(k\) any commutative ring of characteristic
\(p\). No field, reducedness or perfection assumption is imposed. Define
\[
 \JJ_B(k)=k[z_1,\ldots,z_n]/(z_1,\ldots,z_n)^{B+1},\qquad
 I=(z_1,\ldots,z_n)\JJ_B(k).
\]
The number of variables is finite and \(B\ge0\). Frobenius
\(\FF(a)=a^p\) is a ring endomorphism. Its iterates kill every element of
\(I\) once \(p^h>B\), because \(I^{B+1}=0\). This is nilpotence on the
augmentation ideal, not nilpotence or injectivity on the coefficient ring.
More generally, the algebraic lemmas and vanishing theorem in this section hold
in any commutative characteristic-\(p\) ring with an ideal \(I\) satisfying
\(I^{B+1}=0\); \(\JJ_B(k)\) is the application of interest.

\begin{lemma}[Finite Artin--Schreier resolvent]
\label{lem:resolvent}
Let \(q=p^r\), \(r\ge1\), and \(L\in I\). The equation
\(Y-Y^q=L\) has a unique solution \(Y\in I\). For any \(h\ge0\) such that
\(q^h>B\), it is
\begin{equation}
 Y=\sum_{a=0}^{h-1}L^{q^a}. \label{eq:resolvent}
\end{equation}
\end{lemma}
\begin{proof}
Subtracting the \(q\)-th power of the displayed sum telescopes to
\(L-L^{q^h}=L\). If two solutions differ by \(D\in I\), then
\(D=D^q=D^{q^h}=0\). When \(B=0\), the ideal \(I\) is zero and the empty
sum with \(h=0\) is permitted.
\end{proof}

Let \(L=(L_i)_{i\in A}\) be a finite family in \(I\), and let
\(m_i\in\mathbb N\). Read the multiplicities in \(k\); they are fixed by
Frobenius. Define a weighted pairing and its moments by
\[
 b(U,V)=\sum_{i\in A}m_iU_iV_i,\qquad
 M_e=b(L,\FF^eL)=\sum_i m_iL_i^{p^e+1}.
\]
Only the fact \(m_i^p=m_i\) is used; arbitrary weights in \(k\) would not
suffice.

\begin{lemma}[Pairing of two Frobenius iterates]
\label{lem:pairing}
For all \(a,b\ge0\),
\begin{equation}
 b(\FF^a L,\FF^b L)
     =M_{|a-b|}^{p^{\min(a,b)}}. \label{eq:pairing}
\end{equation}
\end{lemma}
\begin{proof}
For \(a\le b\), Frobenius multiplicativity, additivity and the fixed weights give
\[
 \sum_i m_iL_i^{p^a}L_i^{p^b}
 =\left(\sum_i m_iL_iL_i^{p^{b-a}}\right)^{p^a}.
\]
For \(b<a\), use symmetry. No Frobenius inverse or cancellation is involved.
\end{proof}

\needspace{14\baselineskip}
\begin{theorem}[Finite orthogonality certificate]
\label{thm:orthogonality}
Use the above notation and let \(\lambda\in k\).
\begin{enumerate}
\item \textbf{Inert case.} Suppose \(Y_i-Y_i^{p^2}=L_i\), with \(Y_i\in I\).
If \(M_e=0\) for every odd \(e\) with \(p^e+1\le B\), then
\begin{equation}
 \lambda\sum_i m_iY_i^{p+1}=0\quad\text{in }\JJ_B(k). \label{eq:inert}
\end{equation}
\item \textbf{Ramified case.} Suppose \(Y_i-Y_i^p=L_i\), with \(Y_i\in I\).
If \(M_e=0\) for every \(e\ge0\) with \(p^e+1\le B\), then
\begin{equation}
 \lambda\sum_i m_i\left(Y_i^2-\tfrac12L_i^2\right)=0
 \quad\text{in }\JJ_B(k). \label{eq:ramified}
\end{equation}
\end{enumerate}
\end{theorem}
\begin{proof}
Every moment with \(p^e+1>B\) is zero automatically by degree. In the inert
case, Lemma~\ref{lem:resolvent} expresses \(Y\) as a finite sum of even
Frobenius iterates of \(L\), and \(\FF Y\) as a finite sum of odd iterates.
Each cross-pairing is zero by Lemma~\ref{lem:pairing}: its gap is odd. Hence
\(b(Y,\FF Y)=0\), which gives \eqref{eq:inert}.

In the ramified case, \(Y\) is a finite sum of all iterates of \(L\). Every
pairing of such iterates vanishes, including equal-index pairings through
\(M_0\). Thus \(b(Y,Y)=0\) and \(b(L,L)=M_0=0\), giving
\eqref{eq:ramified}. The inverse of two exists because \(p\) is odd. The scalar
\(\lambda\) stays outside the pairing calculation; it need not be Frobenius-fixed.
\end{proof}

\begin{corollary}[Ideal form]
\label{cor:ideal}
In the inert case the raw expression \(\sum_i m_iY_i^{p+1}\) belongs to the
ideal generated by \(M_e\) for odd \(e\) with \(p^e+1\le B\). In the ramified
case \(\sum_i m_i(Y_i^2-L_i^2/2)\) belongs to the ideal generated by all
\(M_e\) with \(p^e+1\le B\).
\end{corollary}
\begin{proof}
Apply Theorem~\ref{thm:orthogonality} after quotienting by the indicated ideal.
Frobenius descends to any quotient ring, and the image of \(I\) is still
nilpotent. Vanishing in this quotient is exactly ideal membership.
\end{proof}

The substantive simplification is the ideal certificate, not just replacing an
infinite sum by a finite one. We never identify all coefficients
\(\lambda_{e,j}\), classify the support of a product of prime-power series, or
put an exact canonical exponent into the theta factorization. The finite double
sum is used only to prove vanishing by orthogonality.

\section{Obtaining the raw certificate from the existing source}
\label{sec:raw}

\subsection{The coefficient ring and the one-curve inputs}

Use \(k=\OO/(\pi)\), with \(\pi\mid p\). This quotient is generally
\emph{nonreduced}. For example, a square root of \(\pi\) in \(\mathbb C_p\)
has a nonzero class whose square is zero. The existing \code{reducePiPrime}
definition and its characteristic-\(p\) instance use this very quotient, not
the residue field. The lemmas above were deliberately stated at this generality.
\src{5}

Before the double-series expansion, the source already provides:
\begin{align}
 &Y,\Xi,W_0\in\OO[[t]],\qquad Y(0)=\Xi(0)=0,\notag\\
 &\theta(\varpi t)=W_0(t)\exp(c\Xi(t)),\label{eq:rawcurve}\\
 &\overline Y-\overline Y^{q}=t,\label{eq:rawparameter}\\
 &\overline\Xi=
 \begin{cases}
   \lambda\overline Y^{p+1},&q=p^2\quad\text{(inert)},\\
   \lambda(\overline Y^2-t^2/2),&q=p\quad\text{(ramified)}.
 \end{cases} \label{eq:rawresidue}
\end{align}
Bars mean reduction modulo \(\pi\). The exact inputs to retain are:
\begin{itemize}
\item \code{exists_scaled_formalParam}, using its functional-equation field
\eqref{eq:rawparameter}; \src{9}
\item \code{exists_rescale_theta_eq_mul_expSubst}; \src{5}
\item \code{exists_seriesToCp_eq_exponentSeries_of_inert} and
\code{exists_seriesToCp_eq_exponentSeries_of_ramified}. \src{5}
\end{itemize}
The ramified theorem asks for the square coefficient in degree two. From the
resolvent, \(\overline Y=t+O(t^p)\), so
\(\coef{t^2}{\overline Y^2}=1\) immediately. No full support computation is
needed for that premise.

At inert primes, \(\lambda\) is a unit. The existing theorem
\code{norm_coeff_exponentSeries_eq_one} says that
\(\coef{t^{p+1}}{\Xi}\) is a unit. Since
\(\overline Y=t+O(t^{p^2})\) and \(p\) is odd,
\(\coef{t^{p+1}}{\overline Y^{p+1}}=1\). Equation~\eqref{eq:rawresidue}
therefore identifies the residue of that unit with \(\lambda\).
This extracts the needed unit without computing any other coefficient. \src{5}

\subsection{Finite products and resolution moments}

For an integral resolution, let \(L_i(z)=u_i^\dagger z\), and let \(m_i\) be
its natural-number multiplicities. Retain the literal raw exponent
\begin{equation}
 \Psi(z)=\sum_i m_i\Xi(L_i(z)). \label{eq:rawproduct}
\end{equation}
The finite-product theorem
\code{rescale_thetaProd_eq_mul_expSubst} already proves the exact factorization
\begin{equation}
 F(\varpi z)=W(z)\exp(c\Psi(z)),\qquad
 W,\Psi\in\OO[[z]],\quad\Psi(0)=0. \label{eq:rawfactor}
\end{equation}
For the theta product alone,
\(W=\prod_i W_0(L_i)^{m_i}\). An additional permitted multiseries factor
becomes integral after rescaling and multiplies \(W\). It does not require a
change in the exponent. \src{6}

The crucial retained resolution identity is
\begin{equation}
 M_e(z)=\sum_i m_i L_i(z)^{p^e+1}\equiv\Gamma_e(z)\pmod\pi
 \label{eq:resolution}
\end{equation}
for every relevant partner exponent. This is precisely
\code{piPrimeInt_dvd_sum_linearSeries_sub_coupling}. All multiplicities in this
statement and the theta-product statement are explicitly natural numbers, so the
fixed-weight assumption in Lemma~\ref{lem:pairing} matches the source. \src{6}

The choices \(\varpi,Y,\Xi,\lambda,W_0\) can be made once for the curve and
the prime. The raw \(\Psi\) is allowed to depend on the resolution. The audited
consumers each use one resolution at a time; none requires identical raw
exponents for different resolutions. Thus dropping the canonical normal form
does not introduce a missing uniformity requirement.

\section{Divisibility: the finite contract is sufficient}
\label{sec:upper}

\begin{theorem}[Null-substitution integrality]
\label{thm:integrality}
Let \(K\supset\mathbb Q_p\) be a valued field with valuation ring \(\OO\),
\(v(p)=1\), and \(0\ne\pi\in\OO\) with \(\pi\mid p\). Suppose
\eqref{eq:rawfactor} holds and that the residue of \(\Psi\) has the inert or
ramified form induced by \eqref{eq:rawresidue} and \eqref{eq:rawproduct}. Assume
\begin{equation}
 v(c\pi)\ge\frac1{p-1}. \label{eq:radius}
\end{equation}
Let every component of \(Z(t)\) be an integral formal series with zero constant
term. If the relevant moments \(M_e(Z)\) vanish modulo \(\pi\) through total
degree \(B\), for all relevant \(e\) with \(p^e+1\le B\), then every
coefficient of \(F(\varpi Z)\) of total degree at most \(B\) is integral.
\end{theorem}
\begin{proof}
Zero-constant substitution preserves the augmentation filtration. Work in
\(\JJ_B(k)\) for \(k=\OO/(\pi)\), using \(L_i(Z)\in I\). These series need
not remain linear. Equations~\eqref{eq:rawparameter}--\eqref{eq:rawresidue} and
Theorem~\ref{thm:orthogonality} imply that the degree-\(B\) jet of \(\Psi(Z)\)
is zero modulo \(\pi\). Its jet over \(\OO\) can therefore be written
\(\pi H\), with \(H\in I\JJ_B(\OO)\).

In \(\JJ_B(K)\), the formal exponential is the finite sum
\[
 \exp(c\Psi(Z))=\sum_{a=0}^{B}\frac{(c\pi)^a H^a}{a!}.
\]
For each \(a\), Legendre's bound gives
\[
 v\left(\frac{(c\pi)^a}{a!}\right)
 \ge\frac{a}{p-1}-v_p(a!)\ge0.
\]
Thus every coefficient in the required range is integral. Multiplication by the
integral series \(W(Z)\) preserves the conclusion.
\end{proof}

The source has the exact threshold
\begin{equation}
 v(c\pi)=\frac1{p-1}, \label{eq:exactradius}
\end{equation}
proved as \code{norm_pow_succ_div_mul_piPrimeInt}. At an inert prime,
\(v(\varpi)=1/(p^2-1)\) and \(v(\pi)=1\). At the ramified prime,
\(v(\varpi)=1/(2p-2)\) and \(v(\pi)=1/2\). Both substitutions give
\eqref{eq:exactradius}. There is no extra valuation loss. \src{6}

In \code{norm_jet_minorProd_le_of_isChargeExponent}, Step~2 currently invokes
\code{HasNormalForm.norm_coeff_subst_le_one}. The null-frame construction has
already made \(\Gamma_e(Z)\) divisible by \(\pi\) for
\(e\le\lfloor\log_p\max(B,1)\rfloor\) in the relevant set. This is the source's
natural-number logarithm, including its value zero at \(B=0\). It includes every
\(e\) with \(p^e+1\le B\). Equation~\eqref{eq:resolution} and
Theorem~\ref{thm:integrality} therefore supply exactly the same coefficient
bound at this call site. \src{10}

The chain-rule determinant cost, charge count, Legendre factorial estimate,
\(\gamma_p\), and the error
\(fN^3(1+\lfloor\log_p\max(B,1)\rfloor)\) can all be retained. The proposal
changes neither the choice of frames nor the quantitative statement of this
local bound. This is why it composes with the proposed global improvements
without requiring a new parameter optimization.

\section{Sharpness: retain a separate gradient certificate}
\label{sec:sharp}

\subsection{The finite differential calculation}

Differentiation lowers the degree filtration by one. To determine a derivative
through degree \(D\), use the exponent through degree \(D+1\). In general,
differentiation is not an endomorphism of \(\JJ_D(k)\). In what follows
\(\dd\) differentiates the variables while fixing the coefficient ring \(k\).

At an inert prime, differentiate the parameter equation to obtain
\(\dd Y_i=\dd L_i\), since \(\dd(Y_i^{p^2})=0\). Hence
\[
 \dd\overline\Psi
 =\lambda\sum_i m_iY_i^p\,\dd L_i.
\]
The finite resolvent for \(Y_i\) gives, through degree \(D\),
\begin{equation}
 \dd\overline\Psi
 =\lambda\sum_{\substack{e\text{ odd}\\p^e\le D}}\dd M_e.
 \label{eq:gradient}
\end{equation}
Indeed \(Y_i^p\) is the sum of the odd Frobenius iterates of \(L_i\), while
\(\dd M_e=\sum_i m_iL_i^{p^e}\dd L_i\) for \(e>0\). Every omitted term
has coefficient degree greater than \(D\). This proof uses neither the
two-index expansion nor an interchange of an infinite sum and a derivative.

Using \eqref{eq:resolution} for the original linear forms gives
\begin{equation}
 \partial_{z_y}\overline\Psi
 =\lambda\sum_{\substack{e\text{ odd}\\p^e\le D}}
       \sum_{\operatorname{pr}_e(x)=y}z_x^{p^e}
 \quad\text{through degree }D. \label{eq:signature}
\end{equation}
The first factors of \(\Gamma_e\) differentiate to zero in characteristic
\(p\). A pure monomial \(z_x^{p^e}\) determines \((x,e)\) uniquely, so its
coefficient in \eqref{eq:signature} is exactly \(\lambda\). In particular,
there are no unintended multiplicities from the partner map.

For completeness, at the ramified prime the analogous calculation is
\begin{equation}
 \dd\overline\Psi
 =\lambda\left(\tfrac12\dd M_0+
       2\sum_{\substack{e\ge1\\p^e\le D}}\dd M_e\right)
 \quad\text{through degree }D. \label{eq:ramgradient}
\end{equation}
It follows from \(\dd(Y_i^2-L_i^2/2)=(2Y_i-L_i)\dd L_i\).
The inert coupling-sharpness argument is the consumer that requires the unit
gradient certificate; no new ramified unit assertion is needed.

\subsection{The coefficient witness, without a canonical exponent}

Let \(\mathfrak m\) be the maximal ideal of \(\OO\). Congruence modulo
\(\pi\) implies congruence modulo \(\mathfrak m\), because
\((\pi)\subset\mathfrak m\). For sharpness, use
\eqref{eq:rawfactor} with \(W\) a unit and \(v(c)<0\). At the inert primes
under consideration, \(v(c)=1/(p-1)-1<0\). Unitness of \(W\) is required
here, not merely integrality; in the normalized theta product it follows from
constant term one.

Consider the cancelling product in two sets of variables
\[
 G(z',z'')=F(z'+z'')\prod_{a\ne0}F(\delta_a z'),\qquad
 U(z')=G(z',0),
\]
where the existing cancellation construction supplies integral \(U\) with
\(U(0)=1\). For this discussion \(F=W\exp(c\Psi)\) denotes the already
rescaled function. The cancellation factors can be those of the current source
or the fewer factors supplied by the existing CM-norm proposal.

The logarithmic derivative in a \(z''\)-variable is
\[
 \partial_yG=(cA_y+B_y)G,
 \qquad A_y=(\partial_y\Psi)(z'+z''),\quad B_y\text{ integral}.
\]
If \(g\) is integral, then
\begin{equation}
 \partial_y(gG)=c\,g'G,\qquad
 g'=gA_y+c^{-1}(\partial_yg+gB_y),\qquad
 \overline{g'}=\overline g\,\overline{A_y}\pmod{\mathfrak m}.
 \label{eq:recurrence}
\end{equation}
Here \(c^{-1}\in\mathfrak m\). Repetition gives the product of the first
derivatives modulo \(\mathfrak m\); higher derivatives appear only in terms
killed by this factor. The existing
\code{exists_jet_prod_copySubst_couplingSubst} already proves this for an
\emph{arbitrary integral exponent}. \src{11}

For each selected place \(\tau\), take an odd exponent \(e_\tau\), partner
\(y_\tau=\operatorname{pr}_{e_\tau}(\tau)\), and multiplicity \(j_\tau\).
The existing \code{IsCouplingFamily} condition says that, among selected
variables, \(z_\tau^{p^{e_\tau}}\) is the unique lowest-degree monomial of
\(\partial_{y_\tau}\overline\Psi\). Equation~\eqref{eq:signature} supplies
its unit coefficient. Put
\[
 J=\sum_\tau j_\tau,\qquad
 \nu'=\sum_\tau p^{e_\tau}j_\tau\mathbf e_\tau,\qquad
 \nu''=\sum_\tau j_\tau\mathbf e_{y_\tau}.
\]
After \(J\) derivatives in \(z''\), the reduction of \(g(z',0)\), restricted
to selected variables, is
\(\prod_\tau(\partial_{y_\tau}\overline\Psi)^{j_\tau}\).
Its unique lowest term has exponent \(\nu'\) and unit coefficient.
Multiplication by \(U\), whose constant term is one, preserves that coefficient.
Consequently
\[
 \nu''!\;\coef{(z')^{\nu'}(z'')^{\nu''}}G
       =c^J\,\coef{(z')^{\nu'}}{g(z',0)U(z')},
\]
and the coefficient on the right after \(c^J\) is a unit. Thus
\begin{equation}
 v\!\left(\coef{(z')^{\nu'}(z'')^{\nu''}}G\right)
       =Jv(c)-v(\nu''!). \label{eq:sharpvalue}
\end{equation}
This is the same factorial and valuation identity as
\code{IsCouplingFamily.norm_coeff_pairExponent}. \src{11}
Here \(\nu''!\) means the multiindex factorial
\(\prod_y(\nu''_y)!\), not the factorial \(J!\).

The only specialization to \code{couplingSum} that must be replaced is the
proof of the first-gradient lowest-term assertion and the signatures carrying
it. The generic recurrence, cancellation product, factorial conversion and
least-valuation bookkeeping remain valid. One does not need to take successively
larger jets at every derivative: \eqref{eq:recurrence} uses exact integral
series, and only their first-gradient reductions enter the leading-term test.

\section{Scaling, uniformity, and the new interface}
\label{sec:interface}

\subsection{Two contracts, with a shared exact factorization}

A useful interface stores an existential integral \(\Psi\), integral \(W\),
\(\Psi(0)=0\), and the exact factorization \eqref{eq:rawfactor}. It exposes:
\begin{description}[style=nextline,leftmargin=0pt,labelwidth=0pt,labelsep=0pt]
\item[Null-moment contract.]
For every finite degree bound and zero-constant integral substitution, vanishing
of the relevant coupling jets modulo \(\pi\) forces vanishing of the exponent
jet modulo \(\pi\).
\item[Inert gradient contract.]
Every nonzero residue coefficient of \(\partial_y\Psi\) has monomial
\(z_x^{p^e}\), with odd \(e\) and \(\operatorname{pr}_e(x)=y\); at each
prescribed monomial the coefficient is the residue of a specified unit \(u_e\).
\end{description}
Initially \(u_e=\lambda\) for every odd \(e\). Allowing a family of units
\(u_e\) makes transport under scaling immediate. Upper-bound callers do not
need the unit-gradient field; sharpness callers additionally require \(W\) to
be a unit. These distinctions should be reflected in theorem hypotheses.

\subsection{Ordinary integral scaling}

Replace \(F(z)\) by \(F(\alpha z)\), with \(\alpha\in\OO\), keeping
\(c\) and \(\varpi\) fixed. Then
\[
 W'(z)=W(\alpha z),\qquad \Psi'(z)=\Psi(\alpha z),\qquad
 M'_e=\alpha^{p^e+1}M_e.
\]
The null-moment contract survives even when \(\alpha\) is a nonunit or zero.
The gradient coefficients transform as
\(u'_e=\alpha^{p^e+1}u_e\). They remain units when \(\alpha\) is a unit.
This matches the integral rescaling used in the vacuum/divisibility consumer
and the unit Taylor-scale rescaling used in the sharpness consumer. \src{12}

\subsection{Changing the chosen scale}

If \(\varpi'=\alpha\varpi\) with \(\alpha\in\OO^\times\), then
\(c'=\alpha^{p+1}c\). Take
\[
 W'(z)=W(\alpha z),\qquad
 \Psi'(z)=\alpha^{-(p+1)}\Psi(\alpha z).
\]
This gives the exact factorization with \(\varpi'\), and the gradient units
become \(u'_e=\alpha^{p^e-p}u_e\). Both contracts persist.

One must not claim that the normalized parameter equation is unchanged under
this rescaling. If \(Y'(L)=Y(\alpha L)/\alpha\), then
\[
 Y'-\alpha^{q-1}(Y')^q=L.
\]
Construct the certificate at the canonical scale and transport its consequences,
or explicitly allow this twisted equation. The current
\code{HasNormalForm.exists_of_norm_eq} has no external lexical caller in the
audited snapshot, but its useful covariance is still available in the replacement.

\section{Source changes and realistic savings}
\label{sec:migration}

\subsection{Affected regions}

The table counts physical lines, including comments and blank lines, at the
pinned commit. Paths are relative to \code{Lean/grh/src}. Exact ranges,
declarations, hashes and lexical occurrences are in
\code{audit/source_audit.json}. The count is a proposed replacement boundary,
not an applied patch or an elaborated dependency graph.

\begin{center}
\small
\begin{tabularx}{\linewidth}{@{}X r X@{}}
\toprule
File & Lines & Intended change\\
\midrule
\code{Lib/General/Algebra/}\newline\code{SeriesCongruences.lean}
&199& Remove the specialized prime-power-series support and reindexing section
once its old exponent callers disappear.\\
\code{New/Derivatives/FormalGroup/}\newline\code{Exponent.lean}
&185& Replace the prime-power expansion wrappers and final double-series constructor
with the raw one-curve certificate; retain the raw congruences and unit proof.\\
\code{New/Derivatives/Divisibility/}\newline\code{NormalForm.lean}
&276& Replace \code{couplingSum}, the exact normal-form interface and its transport,
divisibility and gradient machinery by the two contracts.\\
\code{New/Derivatives/FormalGroup/}\newline\code{ThetaProducts.lean}
&92& Replace the sum-to-couplingSum congruence and normal-form constructor.
Keep finite product factorization and the resolution moment identity.\\
\midrule
\textbf{Primary affected regions}&\textbf{752}&Before replacement code and adapters.\\
Optional helper cleanup&84&Continuity, summability and divisibility helpers whose
specialized callers disappear; verify their resolved dependencies before deletion.\\
\bottomrule
\end{tabularx}
\end{center}

The 314-line checked core bundled here is deliberately an independent,
generally reusable development with auxiliary statements and comments. It is
not the total implementation cost of a port. Additional quotient/jet, derivative,
formal-group, and caller adapters are still required, while parts of the existing
generic infrastructure can be reused. It would be misleading to subtract 314
from 752 and call the result a measured net saving.

The defensible claim is substantial simplification of this normalization stage:
one entire exact classification layer becomes unnecessary, and its two consumers
have direct proofs. There is no measured whole-repository percentage reduction
or kernel-time improvement yet. This proposal does not eliminate the major
automorphic, Galois, determinant or analytic parts of the proof.

\subsection{A staged implementation}

\begin{enumerate}
\item \textbf{Introduce the finite algebra.} Port the standalone ring-endomorphism
lemmas. Instantiate Frobenius on the finite jet quotient or use coefficientwise
truncation with an equivalent nilpotence lemma. Natural multiplicities require no
new fixed-weight hypothesis.
\item \textbf{Construct the raw one-curve data.} Combine
\eqref{eq:rawcurve}--\eqref{eq:rawresidue}. Recover the ramified degree-two
premise and the inert unit from the finite low-degree argument. Keep
\code{norm_coeff_exponentSeries_eq_one}.
\item \textbf{Construct resolution certificates.} Use
\eqref{eq:rawproduct}, the exact product theorem and the existing moment congruence.
Prove the null-moment and gradient fields directly, then their scaling lemmas.
\item \textbf{Switch the divisibility caller.} Change Step~2 of the main count
theorem to Theorem~\ref{thm:integrality}. Leave its frame, determinant and numerical
estimates intact. Update the wrappers in \code{Vacuum/Divisibility.lean}.
\item \textbf{Switch the sharpness caller.} Generalize the lowest-term lemmas and
\code{norm_coeff_pairExponent} to an arbitrary exponent with the gradient
certificate. Keep the already generic logarithmic-derivative recurrence. Update
the constructor in \code{Sharpness/Cases.lean}.
\item \textbf{Remove the old dependency path.} Once both branches compile,
delete the old normal-form route and audit now-unused helpers. A temporary
compatibility theorem from the old normal form is useful during migration, but
it cannot remain the proof of the final new constructor if actual simplification
is claimed.
\item \textbf{Replay the final endpoints.} Build the affected modules and the
exported RH/Dirichlet-GRH declarations; inspect their axioms. Only after this
step report net lines, dependency reduction and build-time measurements.
\end{enumerate}

An important retention: the all-coefficients description of
\code{exists_scaled_formalParam} also has a consumer in
\code{Sharpness/OneCurve.lean}. The new argument does not need that stronger
description, but this does \emph{not} authorize deleting it from the repository.
Similarly, generic results on integral power series should be retained unless
their actual elaborated callers are gone.

\section{Validation and limits of the deliverable}
\label{sec:validation}

\subsection{What has been formally checked}

\code{formal/FrobeniusJets.lean} compiles with
Lean \code{4.35.0} and mathlib commit
\code{56a005855ea17e0b0958f73ec801ae140b0a1679}, exactly the versions requested
by the inspected project. It proves:
\begin{itemize}
\item the finite geometric-resolvent identity, including its tail;
\item existence and uniqueness on a nilpotent endomorphism orbit;
\item the weighted pairing identity and finite double-pairing identity;
\item the inert raw-exponent vanishing theorem from odd moments;
\item the ramified corrected raw-exponent vanishing theorem from all moments;
\item the specialization to Frobenius with natural multiplicities.
\end{itemize}
Its printed theorem axioms are precisely
\code{propext}, \code{Classical.choice}, and \code{Quot.sound}.
There are no added axioms or unproved placeholders. The file is abstract over a
commutative ring and an endomorphism, so it neither imports the GRH development
nor silently assumes its conclusions.

The checked file does not yet construct the multivariate jet quotient, prove the
valuation-ring exponential adapter, formalize the gradient calculation, or
replace the source consumers. Those parts are proved mathematically here and
matched to their source interfaces. A standalone checked algebraic component is
not a completed formal refactor. The exact compile command, tool versions,
source hash, successful exit status, and axiom output are included in
\code{formal/validation.txt}.

\subsection{Exact finite-ring regressions}

\code{checks/check_finite_jets.py} runs without third-party packages. It uses
\[
 (\mathbb F_p[\epsilon]/\epsilon^2)[z_1,z_2]/(z_1,z_2)^{B+1},
 \qquad p\in\{3,5,7\},
\]
including cutoff values immediately before and after prime powers. The
coefficient nilpotent \(\epsilon\) tests that no hidden reducedness or
Frobenius-injectivity assumption is present. The common scalar \(\lambda\)
is deliberately allowed not to be Frobenius-fixed.

There are 29 parameter cases, three deterministic randomized trials per case,
and 2,614 assertion groups: 1,392 pair identities, 522 resolvent checks,
174 uniqueness checks, 87 inert and 87 ramified exponent identities,
348 gradient comparisons using the correct input precision, and four negative
controls. All pass. The seed and machine-readable results are in
\code{checks/finite_jet_results.json}. These tests are useful regressions for
indices and coefficients; the universal proofs are the mathematical arguments
and the applicable checked Lean lemmas.

\subsection{Source and proposal audit}

The source audit searched all 803 Lean source files for the relevant declarations
and callers. The exclusion inventory covers 18 proposal directories and 64
textual proposal files, including 28 TeX sources, at the same pinned commit.
The manifest records byte lengths, Git blob hashes and SHA-256 hashes; the
bundled verifier checks them against a supplied checkout. The audit does not
claim independent verification of the existing proposals' proofs or executables.

The global proof was not rebuilt during this work. The mathematical result of
this report is a replacement argument with matched local contracts; the
engineering result is the compiled finite algebraic core and a source-specific
migration plan. A full port is the next step required to turn the proposed
simplification into a smaller checked GRH dependency graph.

\appendix
\section{Why this does not duplicate the existing proposals}
\label{app:inventory}

The inventory below is relative to the pinned snapshot, not older inventories
embedded in previous reports. In particular it includes
\code{discrete_hecke_simplification}, added at that snapshot. It lists what was
treated as already proposed and therefore excluded as an answer to the request.
All entries refer to the source proposal directory \src{13}; the archive's
\code{audit/proposal_inventory.tsv} gives the complete mechanism comparison.

\small
\begin{longtable}{@{}p{0.36\linewidth}p{0.60\linewidth}@{}}
\toprule
Directory & Mechanisms already proposed\\
\midrule
\endfirsthead
\toprule Directory & Mechanisms already proposed\\\midrule
\endhead
\bottomrule\endfoot
\code{centered_theta_lift_}\newline\code{simplification} & Centered covariance, direct ramified-source
spanning, adaptive completion, product-of-contractions variant and a smaller
logarithmic translate budget.\\[5pt]
\code{cm_norm_theta_simplification} & CM-norm cancellation, two or three cancelling
factors, and four-copy local sharpness in place of the previous twelve-copy bound.\\[5pt]
\code{cyclotomic-interpolation-and-theta-lifts} & Mixed-radius torsion grids,
dimension/degree-independent order control, one character for all powers, and
stronger near-spanning from two moments.\\[5pt]
\code{discrete_hecke_simplification} & Signed integer-lattice enumeration instead
of Pieri-cone integration, plus direct saturation/Smith control of exceptional
local elementary divisors.\\[5pt]
\code{edge_hecke_simplification} & Integral orbit traces and coefficient operators
instead of restricted-representation machinery; low-rank/principal-submatrix
control of exceptional Smith factors.\\[5pt]
\code{gram_adjoint_simplification} & Inverse Gram duality for an integral adjoint,
removing the tight-frame construction and improving hook normalization.\\[5pt]
\code{grh-class-count-finite-quotient} & Finite-cokernel/residue invariants replacing
the bounded strong-approximation class-count construction.\\[5pt]
\code{grh_order_tests_2026-11-27} & Norm power tests and one subgroup exclusion
replacing primitive labels, with elementary counting and all full-order consumers
retained.\\[5pt]
\code{newton_torsion_simplification} & One-shell Newton divided differences,
weighted Gauss valuations and harmonic averaging for a common primitive order.\\[5pt]
\code{orbital_theta_exchange} & Normalized orbital sections, algebraic twisted
coinvariants and canonical matching instead of representative choices and
extended finite sums.\\[5pt]
\code{proof_simplifications} & Large-sieve or Fourier auxiliary-prime selection,
far-cutoff Perron, and direct spanning with stronger missing-dimension control.\\[5pt]
\code{rh-simplifications} & Direct spectral rank recurrence, simultaneous torsion
interpolation, quadratic-large-sieve selection and endpoint assembly.\\[5pt]
\code{rh_sparse_theta_proof} & Sparse/dyadic cutoffs, reciprocal-cutoff convexity,
local-average arithmetic criteria and Gaussian-tail twist selection.\\[5pt]
\code{rh_two_order_tests_proposal} & Two surviving order components, weaker tame
order certificates, exact cyclic counting and improved projector complement loss.\\[5pt]
\code{robust_theta_labels} & Three power tests, robust labels and their
prime-power/Galois transfer, replacing analytic totient estimates.\\[5pt]
\code{shorter-theta-lift-proof} & Elementary smooth character blocks or the
quadratic large sieve, direct spanning, fixed-modulus far-cutoff PNT and interpolation.\\[5pt]
\code{signed_theta_gain_simplification} & Signed product-formula normalization,
separation of divisibility loss and colength, and smaller explicit global loss.\\[5pt]
\code{theta-lift-simplifications} & Integral flag-row determinants and compressed
lifts, denser robust labels, and an explicit paraboloid/digit place construction.\\
\end{longtable}
\normalsize

Their stated future directions were also treated as exclusions. These include
joint rank--height optimization, weaker orbital estimates, charging index by new
rank, exterior-power determinant arguments avoiding completion, direct perfect
pairings on small-defect integral spaces, finite-scale/Ehrhart packing, full
Smith profiles, root-by-root multi-block shifts, and sharper local weighted
losses. This report does not offer one of those open directions as its result.

The nearest overlap is the CM-norm proposal. Its report explicitly retains the
formal-group normal form while changing the cancelling product
(\code{report.tex}, lines 194--195 in that directory). Our replacement removes
the normal-form dependency while leaving that improved cancellation available.
The other occurrence of Frobenius in the discrete/edge Hecke proposals concerns
residue-representation packing, not the formal-group exponent. The two uses must
not be conflated merely because both involve characteristic \(p\).

The novelty claim is relative to this repository's proposal corpus. Nilpotent
Artin--Schreier inversion and the Frobenius pairing identity are elementary
algebra; no historical priority claim is made for them.

\section{Failure modes that the replacement must avoid}
\label{app:failures}

\sectionlead{A gradient alone does not prove integrality.}
In characteristic \(p\), \(\dd(z^p)=0\), although \(z^p\ne0\) in a large
enough jet algebra. More concretely, take \(p=3\), \(\pi=3\), and
\(c=1/\sqrt3\) in \(\mathbb C_3\). Then
\(v(c\pi)=1/2=1/(p-1)\), but the coefficient of \(z^3\) in
\(\exp(cz^3)\) is \(c\), which is nonintegral. Thus the null-moment
certificate cannot be replaced by the derivative certificate.

\sectionlead{The ramified zero-gap moment is indispensable.}
For \(p=3\) and \(B=2\), take \(L=z\). Then \(Y=z\) in the jet algebra,
and \(Y^2-L^2/2=z^2/2\ne0\), while every positive-gap moment vanishes by
degree. Omitting \(M_0\) would make the ramified statement false.

\sectionlead{Differentiation needs one more degree of input.}
In characteristic three, \(z^4=0\) in \(\JJ_3\), while its derivative
computed before truncation is \(z^3\ne0\) in \(\JJ_3\). The legitimate map
is from input precision \(D+1\) to derivative precision \(D\).

\sectionlead{Weights must be fixed by Frobenius.}
In \(k=\mathbb F_3[\epsilon]/\epsilon^2\), take weight \(w=1+\epsilon\)
and \(L=z\). The pairing with indices \(1,2\) is
\((1+\epsilon)z^{12}\), but the cube of the gap-one moment is \(z^{12}\).
The natural weights in the source satisfy the required condition; an arbitrary
weighted generalization would not.

\sectionlead{Zero constant terms are a real hypothesis.}
Frobenius is not nilpotent on constants. The finite resolvent formula and the
moment degree cutoff rely on \(L_i\) and every substituted variable lying in
the augmentation ideal. This requirement cannot be dropped during transport.

\sectionlead{An unrelated determinant shortcut also needs more input.}
One tempting global shortcut is to choose a different twist for each eigenline
and multiply the best local diagonal gains. Separate integral lattices do not
justify this. For example, over \(\mathbb Q_p\), set
\[
 U_0=\begin{pmatrix}p&1\\1&p\end{pmatrix},\quad
 H_1=\operatorname{diag}(p,1),\quad H_2=\operatorname{diag}(1,p),\quad
 G_1=\operatorname{diag}(1,0),\quad G_2=\operatorname{diag}(0,1).
\]
The maps \(G_1H_1^{-1}U_0\) and \(G_2H_2^{-1}U_0\) preserve respectively
\(\langle e_1,pe_2\rangle\) and \(\langle pe_1,e_2\rangle\), with selected
row multipliers equal to one. Nevertheless
\[
 \frac{\det U_0}{(H_1)_{11}(H_2)_{22}}=\frac{p^2-1}{p^2}
\]
is not integral. A common integral model would be additional mathematical input.
This route is therefore not offered as a proved simplification of the source's
twisted-lattice comparison.

\section{Source references and reproducibility}
\label{app:sources}

All repository links below are pinned to
\code{9050b04000647d588d0c78d945180abcd3b0bcf7}. References identify the
source facts being reused; the new algebraic arguments are proved in this report.
Exact file hashes are in \code{audit/source_manifest.json}.

\begin{description}[leftmargin=2.6em,style=nextline]
\item[S1]\label{src:1}
\sourcelink{Lean/grh/src/Main.lean}{Exported theorems, \code{Main.lean}},
\sourcelink{Lean/grh/src/Statements.lean}{\code{Statements.lean}}, and
\sourcelink{Lean/grh/src/New/Main.lean}{the functional-equation bridge}.
\item[S2]\label{src:2}
\sourcelink{docs/guide/theta_lifts_and_dirichlet_zeros.tex}{Human-readable guide,
\code{theta_lifts_and_dirichlet_zeros.tex}}. In particular, the prime-remainder
definition, three volume estimates, local prime-power calculation, sharpness,
and final comparison sections.
\item[S3]\label{src:3}
\sourcelink{Lean/grh/src/New/Reduction/OneSidedBound.lean}{\code{New/Reduction/OneSidedBound.lean}}:
\code{eventually_primeRemainder_le_rpow} and \code{grh}.
\item[S4]\label{src:4}
\sourcelink{Lean/grh/src/New/Reduction/MainInequality.lean}{\code{New/Reduction/MainInequality.lean}}:
\code{star_inequality}, including the definitions of the size and index allowances.
\item[S5]\label{src:5}
\sourcelink{Lean/grh/src/New/Derivatives/FormalGroup/Exponent.lean}{\code{New/Derivatives/FormalGroup/Exponent.lean}}:
raw factorization at line 240; inert/ramified raw congruences at 420/446;
reduction map at 495; characteristic instance at 510; unit coefficient at 566;
double-series construction at 624, 663 and 702.
\item[S6]\label{src:6}
\sourcelink{Lean/grh/src/New/Derivatives/FormalGroup/ThetaProducts.lean}{\code{New/Derivatives/FormalGroup/ThetaProducts.lean}}:
integrality radius at 114; resolution moment congruence at 244; old sum-to-normal-form
bridge at 301; exact raw finite-product factorization at 405; old constructor at 443.
\item[S7]\label{src:7}
\sourcelink{Lean/grh/src/New/Derivatives/Divisibility/NormalForm.lean}{\code{New/Derivatives/Divisibility/NormalForm.lean}}:
\code{couplingSum} at 43; \code{HasNormalForm} at 50; generic finite exponential
coefficient bound at 131; null-substitution consumer at 262; derivative-support
lemmas at 355--417.
\item[S8]\label{src:8}
\sourcelink{Lean/grh/src/Lib/General/Algebra/SeriesCongruences.lean}{\code{Lib/General/Algebra/SeriesCongruences.lean}}:
specialized prime-power-series section, lines 36--234, and generic congruence helpers.
\item[S9]\label{src:9}
\sourcelink{Lean/grh/src/Lib/CM/FormalGroup/FormalParameter.lean}{\code{Lib/CM/FormalGroup/FormalParameter.lean}}:
\code{exists_scaled_formalParam}, line 224. The coefficient-description consumer
that must be retained is in
\sourcelink{Lean/grh/src/New/Derivatives/Sharpness/OneCurve.lean}{\code{Sharpness/OneCurve.lean}},
line 298.
\item[S10]\label{src:10}
\sourcelink{Lean/grh/src/New/Derivatives/Divisibility/Count.lean}{\code{New/Derivatives/Divisibility/Count.lean}}:
\code{norm_jet_minorProd_le_of_isChargeExponent}, line 585, especially its Step~2
call at line 620 and the wrappers at 662 and 707.
\item[S11]\label{src:11}
\sourcelink{Lean/grh/src/New/Derivatives/Sharpness/Couplings.lean}{\code{New/Derivatives/Sharpness/Couplings.lean}}:
generic logarithmic derivative and recurrence at 108--165; coupling lowest-term
lemma at 201; generic exponent theorem at 328; coefficient norm identity at 358.
\item[S12]\label{src:12}
\sourcelink{Lean/grh/src/New/Derivatives/Vacuum/Divisibility.lean}{\code{New/Derivatives/Vacuum/Divisibility.lean}},
especially its constructor at 170, and
\sourcelink{Lean/grh/src/New/Derivatives/Sharpness/Cases.lean}{\code{New/Derivatives/Sharpness/Cases.lean}},
especially \code{exists_hasNormalForm_resolutionForms} at 245.
\item[S13]\label{src:13}
\href{\repo/tree/\commit/Analysis/GeneralizedRiemannHypothesis/docs/proposals}{All 18 existing proposal directories}.
The manifest includes all 64 textual files used in the exclusion review.
\end{description}

\sectionlead{Reproduce the artifacts.}
From the archive root, run
\begin{quote}\small\ttfamily
python3 checks/check\_finite\_jets.py\par
python3 checks/verify\_source\_snapshot.py /path/to/ProveIt\par
latexmk -pdf -interaction=nonstopmode grh\_frobenius\_jet\_simplification.tex
\end{quote}
The formal subdirectory contains a portable lake project and its own setup
instructions. A clean formal run may need to obtain the pinned Lean toolchain,
mathlib and its cache. The supplied validation log records the actual successful
run against an existing cache at those exact versions. The report PDF is already
compiled, so reading it requires no build environment.

\end{document}
